\section{Связность графов. Алгоритмы обхода графов}
\label{sec:graph-connectivity}

\subsection{Компоненты связности}

В неориентированном графе две вершины связаны, если существует соединяющий их
путь. Это отношение эквивалентности: рефлексивность обеспечивается путём длины
ноль, симметричность --- обращением пути, транзитивность --- склейкой путей.
Следовательно, множество вершин разбивается на компоненты связности.

Граф связен, если имеет одну компоненту. Эквивалентно, связный граф содержит
остовное дерево. Поэтому для связного графа с $n$ вершинами
\[
    m\ge n-1.
\]

Ребро называется мостом, если его удаление увеличивает число компонент.
Вершина называется точкой сочленения, если после её удаления вместе с
инцидентными рёбрами число компонент увеличивается.

\textbf{Критерий моста.} Ребро связного неориентированного графа является мостом
тогда и только тогда, когда не лежит ни в одном цикле. Если ребро лежит в цикле,
между его концами остаётся обходной путь; если не лежит, альтернативного пути
нет.

\subsection{Вершинная и рёберная связность}

Вершинная связность $\kappa(G)$ --- минимальное число вершин, удаление которых
делает граф несвязным или тривиальным. Рёберная связность $\lambda(G)$ ---
минимальное число рёбер, удаление которых нарушает связность.

Для нетривиального графа
\[
    \kappa(G)\le \lambda(G)\le \delta(G),
\]
где $\delta(G)$ --- минимальная степень. Правая оценка очевидна: удалив все
рёбра, инцидентные вершине минимальной степени, можно её изолировать. Левая
оценка отражает то, что разделение графа вершинами в общем случае не сложнее,
чем рёбрами.

\subsection{Теорема Менгера}

$u$--$v$-пути называются внутренне вершинно непересекающимися, если их общими
вершинами могут быть только $u$ и $v$. Рёберно непересекающиеся пути не имеют
общих рёбер.

\textbf{Теорема Менгера, вершинная форма.}
Для двух несмежных вершин $u,v$ минимальное число внутренних вершин, удаление
которых разделяет $u$ и $v$, равно максимальному числу попарно внутренне
вершинно непересекающихся $u$--$v$-путей.

\textbf{Рёберная форма.}
Минимальное число рёбер, удаление которых разделяет $u$ и $v$, равно
максимальному числу попарно рёберно непересекающихся $u$--$v$-путей.

Глобальное следствие: граф с не менее чем $k+1$ вершинами $k$-вершинно связен
тогда и только тогда, когда между любыми двумя различными вершинами существует
не менее $k$ внутренне вершинно непересекающихся путей.

Теорема Менгера тесно связана с max-flow min-cut. Для рёберной версии достаточно
положить единичные ёмкости на рёбра; для вершинной версии используют расщепление
каждой вершины $v$ на $v_{in}$ и $v_{out}$ с дугой единичной ёмкости между ними.

\subsection{Поиск в ширину (BFS)}

BFS обходит граф слоями по расстоянию от стартовой вершины $s$. Используется
очередь.

Инициализация:
\[
    dist[s]=0,
\]
остальные расстояния бесконечны. При извлечении вершины $v$ все ещё не открытые
соседи $u$ получают
\[
    dist[u]=dist[v]+1,
    \qquad parent[u]=v.
\]

При списках смежности каждая вершина помещается в очередь не более одного раза,
а каждое ребро просматривается постоянное число раз, поэтому
\[
    T=O(n+m),
    \qquad
    M_{extra}=O(n).
\]

\textbf{Теорема о корректности BFS.}
После завершения обхода
\[
    dist[v]=d(s,v)
\]
для каждой достижимой вершины $v$ в невзвешенном графе.

Идея доказательства: BFS сначала полностью обрабатывает слой расстояния $k$, и
только затем вершины расстояния $k+1$. Если бы вершина получила расстояние больше
кратчайшего, предпоследняя вершина некоторого кратчайшего пути была бы открыта
раньше и обнаружила бы её с правильной меткой.

Массив $parent$ задаёт дерево кратчайших путей от $s$.

\subsection{Поиск в глубину (DFS)}

DFS идёт по неоткрытому соседу настолько глубоко, насколько возможно, после чего
возвращается назад. Реализация --- рекурсивный вызов или явный стек. При списках
смежности
\[
    T=O(n+m).
\]
Для несвязного графа внешний цикл запускает DFS из каждой непосещённой вершины;
получается лес DFS.

Полезно хранить времена
\[
    tin[v],\qquad tout[v].
\]
Для дерева DFS вершина $u$ является предком $v$ тогда и только тогда, когда
\[
    tin[u]\le tin[v]<tout[v]\le tout[u].
\]

В ориентированном графе DFS позволяет классифицировать дуги на древесные,
обратные, прямые и перекрёстные. Наличие обратной дуги в текущего предка является
признаком ориентированного цикла.

\subsection{Мосты и точки сочленения за один DFS}

Для вершины $v$ вводят
\[
 low[v]=\min\Bigl(
 tin[v],
 tin[p]\text{ по обратным рёбрам }(v,p),
 low[u]\text{ по детям }u
 \Bigr).
\]
Величина $low[v]$ --- минимальное время входа в вершину, достижимую из поддерева
$v$ одним подъёмом по обратному ребру (или без него).

Для ребра дерева DFS $(v,u)$:
\[
    (v,u)\text{ --- мост}
    \quad\Longleftrightarrow\quad
    low[u]>tin[v].
\]
Если $low[u]\le tin[v]$, из поддерева $u$ существует альтернативный выход к
$v$ или его предку, значит ребро лежит в цикле.

Для некорневой вершины $v$:
\[
 v\text{ --- точка сочленения}
 \quad\Longleftrightarrow\quad
 \exists u\text{ --- ребёнок }v:\ low[u]\ge tin[v].
\]
Корень DFS является точкой сочленения тогда и только тогда, когда у него не менее
двух детей в дереве DFS. Все мосты и точки сочленения находятся за
\[
    O(n+m).
\]

\subsection{Сильная связность ориентированного графа}

Орграф слабо связен, если после удаления ориентации получается связный граф.
Орграф сильно связен, если для любых $u,v$ существуют пути
\[
    u\leadsto v,
    \qquad
    v\leadsto u.
\]
Компонента сильной связности (КСС) --- максимальное множество попарно взаимно
достижимых вершин.

После стягивания каждой КСС в одну вершину получается \emph{конденсация} графа.
Она всегда ациклична: ориентированный цикл между различными компонентами означал
бы взаимную достижимость и необходимость объединить их.

\textbf{Алгоритм Косарайю:}
\begin{enumerate}
 \item DFS по $G$, вершины записываются по времени выхода;
 \item построить транспонированный граф $G^T$;
 \item запускать DFS по $G^T$ в порядке убывания времени выхода первого обхода.
\end{enumerate}
Каждый запуск третьего шага выделяет одну КСС. Сложность
\[
    O(n+m).
\]
Алгоритм Тарьяна также находит КСС одним DFS, используя стек и значения $low$.

\subsection{Топологическая сортировка}

Топологический порядок орграфа --- линейный порядок вершин, в котором для каждой
дуги $(u,v)$ вершина $u$ стоит раньше $v$.

\textbf{Теорема.} Топологический порядок существует тогда и только тогда, когда
граф является DAG, то есть не содержит ориентированных циклов.

Доказательство необходимости очевидно: цикл потребовал бы
$v_1<v_2<\cdots<v_k<v_1$. Для достаточности можно использовать DFS: вершина
добавляется в список после обработки всех исходящих дуг, а обратный порядок
выхода даёт топологическую сортировку.

Второй способ --- алгоритм Кана: многократно выбирать вершину с нулевой
полустепенью захода, выводить её и удалять исходящие дуги. Если обработаны не все
вершины, оставшаяся часть содержит цикл. Сложность обоих способов
\[
    O(n+m).
\]

\subsection{Что важно сказать на экзамене}

BFS и DFS работают за $O(n+m)$ при списках смежности. BFS доказываемо находит
кратчайшие пути в невзвешенном графе; DFS раскрывает структуру графа. С помощью
$tin/low$ за один DFS находятся мосты и точки сочленения. Для теории связности
центральна теорема Менгера, а для орграфов --- КСС и тот факт, что конденсация
является DAG.

\newpage
